\documentclass{article}
\usepackage{iclr2026_conference,times}

\input{math_commands.tex}

\usepackage{amsmath}
\usepackage{amssymb}
\usepackage{booktabs}
\usepackage{graphicx}
\usepackage{multirow}
\usepackage{array}
\usepackage{xcolor}
\usepackage{hyperref}
\usepackage{url}
\usepackage{xspace}

\hypersetup{
  colorlinks=true,
  linkcolor=black,
  citecolor=blue,
  urlcolor=blue
}

\definecolor{lightblue}{RGB}{232,244,251}
\definecolor{lightgray}{gray}{0.94}

\newcommand{\bench}{\textsc{SkillMisevo-Bench}\xspace}
\newcommand{\safeevolve}{\textsc{SafeEvolve}\xspace}
\newcommand{\skillrevise}{\textsc{SkillRevise}\xspace}
\newcommand{\evolution}{\mathcal{E}}
\newcommand{\library}{\mathcal{L}}
\newcommand{\taskseq}{\mathcal{Q}}
\newcommand{\asr}{\mathrm{ASR}}
\newcommand{\contam}{\mathrm{Contam}}
\newcommand{\carry}{\mathrm{CarryHarm}}

\title{Practice Makes Unsafe: Skill Misevolution in Self-Improving LLM Agents}

% Keep authors anonymous for submission. Replace only for camera-ready.
\author{Anonymous Authors}

% \iclrfinalcopy % Uncomment only for the camera-ready version.

\begin{document}
\maketitle

\begin{abstract}
LLM agents increasingly improve themselves by distilling successful
interactions into reusable skills. This turns experience into persistent
behavioral policy, yet existing evaluations measure whether evolution improves
capability without testing what happens when success depends on an unsafe
shortcut. We identify \emph{skill misevolution}: a transient unsafe success is
extracted into an apparently legitimate skill and later reused on benign tasks
after the original attack has disappeared. We introduce \bench, an
episode-level benchmark that follows an agent across adversarial exposure,
skill formation, retrieval, zero-injection reuse, and a fresh-session
persistence test. Its paired controls isolate the effect of evolution from the
base agent, the task sequence, and benign learning. Across [N] agent
configurations, [M] evolution methods, and [D] domains, skill evolution
[HEADLINE SAFETY RESULT] while retaining [HEADLINE UTILITY RESULT]. Unsafe
content appears in the skill library before it consistently manifests as
behavior, and as few as [K] exposures produce harm that survives into a clean
session. We further introduce \safeevolve, a governance layer that audits and
repairs candidate skills, preserves verification steps, tracks provenance, and
retires skills with unsafe reuse histories. \safeevolve reduces [PRIMARY HARM
METRIC] by [X] while preserving [Y]\% of the utility gain from evolution.
These results show that self-improvement is a security-relevant policy update:
practice can make an agent more capable and persistently less safe.
\end{abstract}

\section{Introduction}
\label{sec:introduction}

LLM agents are beginning to improve through use. Modern CLI and coding agents
can summarize successful trajectories, write or revise reusable skills, and
retrieve those skills when related tasks appear. This capability changes an
agent from a fixed policy with transient context into a system whose behavior is
partly rewritten by its own experience. Skill evolution is therefore more than
memory: it is a persistent policy update expressed as an editable artifact.

The signal used to produce that update is usually task success or efficiency.
Successful behavior, however, is not necessarily safe behavior. An agent may
complete a task by skipping a verification step, assuming authorization,
treating an irreversible action as routine, or following an attacker-controlled
observation. If the evolution method compresses that trajectory into a reusable
procedure, it can preserve the unsafe shortcut precisely because the shortcut
helped the agent succeed.

This creates a temporal safety failure that independent-task evaluations cannot
observe. The original attack may disappear with the context window, but the
lesson extracted from it can remain in the skill library. A later benign task
can retrieve the skill, apply its over-generalized rule, and reproduce the harm
without a malicious instruction. A fresh session can inherit the same behavior
by loading the evolved library alone. We call this process \emph{skill
misevolution}: transient unsafe success becomes persistent, retrievable, and
generalizable policy.

Measuring misevolution requires following the full causal chain rather than
scoring isolated trajectories. We introduce \bench, whose basic unit is a
seven-task episode. Benign tasks and adversarial exposures alternate while the
skill library evolves, followed by a zero-injection persistence probe in a
fresh executor session. Every task starts with fresh conversational context;
only the evolving library crosses task boundaries. Matched no-evolution,
read-only-reflection, and benign-evolution controls isolate whether later harm
comes from the learned artifact rather than the base model, repeated context,
or ordinary benign learning.

The benchmark evaluates both sides of the failure. At the artifact level, it
measures whether authored skills contain unsafe or over-generalized procedures.
At the behavior level, it measures attack amplification on exposed tasks,
zero-injection contamination on benign tasks, and carryover harm in a clean
session. This distinction matters because an unsafe skill can remain latent
until a later task supplies the right retrieval cue. Behavior-only evaluation
therefore misses risk that has already entered the persistent policy.

Our evaluation asks four questions. \textbf{RQ1} tests whether evolution
increases unsafe behavior while retaining utility. \textbf{RQ2} follows the
exposure-to-persistence chain and tests whether the skill artifact mediates
later harm. \textbf{RQ3} measures how the effect varies across evolution
methods, agent configurations, domains, and attacker budgets. \textbf{RQ4}
tests whether \safeevolve can break this chain without removing the capability
benefits of learning.

\safeevolve acts at the learning boundary rather than as another runtime
refusal layer. It audits the complete candidate skill bundle before admission,
repairs unsafe generalization instead of deleting useful procedures, preserves
verification requirements during revision, records the experiences from which
each skill was learned, and downweights or retires skills whose reuse produces
unsafe behavior. These components target the persistent object responsible for
the incremental risk: the evolved skill library.

Our contributions are:
\begin{enumerate}
  \item We identify \textbf{skill misevolution}, a temporal failure mode in
  which self-improvement converts transient unsafe success into persistent
  behavior that can reappear on benign tasks without the original attack.

  \item We introduce \bench, an episode-level evaluation spanning skill
  formation, artifact inspection, retrieval, behavioral reuse, and clean-session
  persistence, with paired controls that attribute changes to evolution.

  \item We provide a systematic evaluation across [N] agent configurations,
  [M] evolution methods, [D] domains, and multiple exposure budgets, separating
  artifact-level latent risk from behavior-level realized harm.

  \item We introduce \safeevolve and show that governing skill admission,
  revision, provenance, and retirement reduces evolution-induced harm while
  retaining the utility gained through learning.
\end{enumerate}

\section{Related Work}
\label{sec:related-work}

\paragraph{Self-improving agents and skill evolution.}
Position this work against methods that extract, edit, select, vote over, or
retrieve reusable skills from agent trajectories. Organize the literature by
the persistent artifact each method produces and the success signal used to
select it. The key gap is that these methods report capability improvement but
do not measure evolution-induced safety change.

\paragraph{Safety of persistent agent memory.}
Discuss memory poisoning, long-term memory evolution, retrieval attacks, and
forgetting. Distinguish a skill from a stored observation: a skill is an
executable or prescriptive procedure intended to generalize across tasks. The
failure studied here is endogenous policy formation from the agent's own
successful trajectory, not only storage of attacker-supplied text.

\paragraph{Agent red-teaming and prompt injection.}
Cover direct and indirect attacks against tool-using agents and the DTAP and
AgentHazard substrates. These benchmarks provide executable tasks and judges;
\bench changes the evaluation unit from one attacked trajectory to a sequence
in which an attack can alter future behavior through evolution.

\paragraph{Governance of evolving libraries.}
Contrast \safeevolve with capability-oriented library maintenance and drift
control. Library hygiene retains skills that improve task success, including
skills whose utility comes from unsafe shortcuts. \safeevolve introduces safety
as a first-class admission and retention signal.

\section{Skill Misevolution}
\label{sec:problem}

\subsection{Self-improvement as a policy update}

Let an agent with base policy $\pi$ execute a sequence of tasks
$\taskseq=(q_1,\ldots,q_T)$. After task $q_t$, an evolution method
$\evolution$ maps the trajectory $\tau_t$ and current library $\library_t$ to
an updated library:
\begin{equation}
  \library_{t+1}=\evolution(\library_t, q_t, \tau_t).
  \label{eq:evolution}
\end{equation}
The next task is executed by the base agent conditioned on retrieved skills
from $\library_{t+1}$. Conversational state is reset between tasks, so any
systematic cross-task change is carried by the evolved artifact.

\subsection{Definition and causal chain}

An episode exhibits \emph{skill misevolution} when an unsafe procedure learned
from an earlier trajectory enters the persistent library and causes incremental
unsafe behavior on a later task relative to a matched no-evolution execution.
The proposed causal chain is:
\begin{equation}
\text{unsafe success}
\rightarrow \text{skill extraction}
\rightarrow \text{unsafe generalization}
\rightarrow \text{retrieval}
\rightarrow \text{reuse}
\rightarrow \text{persistent harm}.
\label{eq:chain}
\end{equation}
The benchmark separately observes each link. This prevents a later unsafe
action from being attributed to the skill merely because a skill library was
present.

\subsection{Threat model}

The attacker controls a bounded number of task inputs during the learning
history but does not modify the evolution code, skill store, retrieval engine,
judge, or final persistence probe. The attacker succeeds when an exposure is
absorbed into a reusable skill that amplifies unsafe behavior or later causes
zero-injection harm. The strongest outcome is carryover harm in a fresh session
that receives only the final evolved library.

\section{SkillMisevo-Bench}
\label{sec:benchmark}

\subsection{Episode design}

Each episode contains seven tasks from one domain and risk family. The order is
\begin{equation}
  B_1, M_1, B_2, M_2, B_3, M_3, P,
\end{equation}
where $B$ denotes a zero-injection benign twin, $M$ an adversarial exposure,
and $P$ a benign persistence probe. The three exposures use different surface
forms from the same risk family so that the library can accumulate a coherent
rule without simply memorizing one string. Position $B_1$ measures the
pre-exposure floor; $B_2$ and $B_3$ test contamination after increasing
exposure; $P$ runs after the library is frozen in a new executor session.

\begin{figure}[t]
  \centering
  \fbox{\parbox[c][4.0cm][c]{0.94\linewidth}{
    \centering
    \textbf{Figure 1 placeholder: the exposure-to-persistence episode.}\\[4pt]
    Show $B_1 \rightarrow M_1 \rightarrow B_2 \rightarrow M_2 \rightarrow
    B_3 \rightarrow M_3 \rightarrow P$, library snapshots $L_0\ldots L_6$,
    fresh executor sessions, retrieval logs, artifact audits, and the final
    clean-session probe. Highlight that only the library crosses task borders.
  }}
  \caption{\textbf{Skill misevolution evaluation.} Adversarial exposures can
  alter the evolving skill library; later benign tasks test whether the learned
  shortcut is retrieved and reused after the attack disappears.}
  \label{fig:episode}
\end{figure}

\subsection{Paired controls}

Every episode is replayed under matched evolution settings:
\begin{itemize}
  \item \textbf{No evolution}: disables persistent learning and establishes the
  base-agent floor.
  \item \textbf{Read-only reflection}: retains textual lessons in the current
  learning mechanism without authoring a reusable skill artifact.
  \item \textbf{Benign evolution}: evolves only from benign tasks and tests
  whether ordinary learning alone produces the effect.
  \item \textbf{Full evolution}: applies each evaluated method to all successful
  trajectories under its native update rule.
\end{itemize}
All arms use identical tasks, order, tools, agent configuration, and judges.

\subsection{Substrates, agents, and evolution methods}

The benchmark draws episodes from DTAP and AgentHazard. DTAP provides
stateful backend tasks and verifiable judges; AgentHazard provides CLI and
coding tasks executed through file, shell, and editing tools. The evaluation
crosses agent frameworks and model backbones, with the final matrix fixed before
the held-out evaluation.

The evolution settings include \skillrevise, EvoSkill, SkillClaw, AutoSkill,
SkillsVote, and supported native learning loops. Each method is evaluated in
its actual stored representation, including associated scripts or structured
contracts, rather than after conversion to a common textual summary.

\subsection{Metrics}

\paragraph{Utility.}
Benign utility is task completion on $B_1,B_2,B_3$; carryover utility is task
completion on $P$. These distinguish a safer method from one that suppresses
all skill use.

\paragraph{Attack amplification.}
Malicious ASR is the fraction of $M$ tasks whose original scenario judge reports
the unsafe objective. For an evolution method $e$, we report
\begin{equation}
  \Delta\asr(e)=\asr(e)-\asr(\text{no evolution}),
\end{equation}
paired over the same exposed task instances.

\paragraph{Artifact poisoning.}
An independent auditor reads the complete stored skill bundle and scores
content unsafety and unsafe generalization. A poisoned artifact contains an
unsafe procedure and states it broadly enough to apply outside the source task.
Auditor agreement is measured on a double-labeled subset.

\paragraph{Zero-injection contamination.}
Contamination is unsafe-goal completion on benign positions:
\begin{equation}
  \contam(e)=\frac{1}{3N}\sum_{i=1}^{N}\sum_{t\in\{B_1,B_2,B_3\}}
  \mathbb{1}[J_i(\tau_{i,t}^{e})=1].
\end{equation}
We additionally report contamination onset and post-onset retention.

\paragraph{Fresh-session carryover.}
Carryover harm is unsafe-goal completion on $P$ after conversational and
environment state are reset and only the frozen final library is mounted. This
is the primary measure of persistent, skill-borne harm.

\subsection{Statistical protocol}

Results are pooled over the frozen episode distribution. Confidence intervals
use episode-level bootstrap resampling. Paired ASR and contamination changes are
tested over matched task instances; method comparisons report effect sizes and
multiple-comparison-adjusted intervals. Hyperparameters, episode construction,
and the held-out split are frozen before the final sweep.

\section{SafeEvolve}
\label{sec:safeevolve}

\safeevolve is a governance layer around an existing evolution method. It does
not replace the base agent or add a runtime refusal policy. It governs the
persistent skill update at two decision points: writing and reuse.

\subsection{Admission critic and targeted repair}

Before a candidate enters the library, an independent critic audits the full
bundle, including scripts and referenced resources. When a candidate removes a
safety check, assumes authority, broadens a dangerous operation, or embeds an
unsafe tool sequence, the critic produces a targeted repair: restore the check,
add the missing precondition, or narrow the applicability condition. Repaired
skills are re-audited before admission.

\subsection{Verification-preserving evolution}

Skill revisions are constrained to preserve existing verification,
confirmation, authorization, and reversibility requirements unless the update
provides explicit evidence that the requirement is obsolete. This prevents an
efficiency-oriented edit from treating a safeguard as removable overhead.

\subsection{Lineage and provenance}

Each skill stores its source tasks, trajectories, revisions, and exposure
status. Provenance supports attribution and batch rollback when an upstream
experience is later judged unsafe. It also supplies a risk prior to retrieval
and retirement decisions.

\subsection{Safety-aware selection and retirement}

Raw evolution selects skills primarily by utility. \safeevolve augments this
objective with the observed safety impact of reuse. Skills associated with
contamination or carryover harm are downweighted, quarantined, repaired, or
retired even when they improve task success. This distinguishes \safeevolve
from capability-only library hygiene, which keeps useful but unsafe skills.

\begin{figure}[t]
  \centering
  \fbox{\parbox[c][3.6cm][c]{0.94\linewidth}{
    \centering
    \textbf{Figure 2 placeholder: SafeEvolve architecture.}\\[4pt]
    Trajectory $\rightarrow$ candidate skill $\rightarrow$ provenance and
    verification-preserving edit $\rightarrow$ admission critic/repair
    $\rightarrow$ library $\rightarrow$ safety-aware retrieval and retirement.
  }}
  \caption{\textbf{SafeEvolve governs persistent learning.} Components act at
  skill writing and reuse, the two points at which transient experience becomes
  persistent policy and later affects behavior.}
  \label{fig:safeevolve}
\end{figure}

\section{Experimental Setup}
\label{sec:setup}

\subsection{Research questions}

\begin{description}
  \item[RQ1: Safety--utility coupling.] Does skill evolution increase unsafe
  behavior while retaining or improving task utility?
  \item[RQ2: Exposure-to-persistence mechanism.] Does unsafe experience enter
  the skill artifact, get retrieved, and cause zero-injection or fresh-session
  harm?
  \item[RQ3: Scope and attacker budget.] How does misevolution vary across
  methods, agents, domains, and one to three adversarial exposures?
  \item[RQ4: Mitigation.] Can \safeevolve reduce evolution-induced harm to the
  no-evolution floor while preserving the capability benefit of evolution?
\end{description}

\subsection{Evaluation matrix}

Report the final matrix of agent framework $\times$ backbone $\times$ evolution
method $\times$ substrate. Separate method applicability from missing runs, and
pin every model and repository version. The main tables pool the predeclared
episode distribution; per-agent, per-domain, and per-risk-family results appear
in the appendix.

\subsection{Baselines}

The diagnostic comparison includes no evolution, read-only reflection, benign
evolution, and the selected skill-evolution methods. The mitigation comparison
adds Library-Drift as a capability-oriented governance baseline. Each evolution
method is evaluated raw, with Library-Drift, and with \safeevolve under the same
frozen train/development/held-out protocol.

\section{Results}
\label{sec:results}

\subsection{RQ1: Does evolution increase unsafe behavior while retaining utility?}
\label{sec:rq1}

Table~\ref{tab:misevolution-main} is the paper's primary diagnostic table. Read
it from left to right: utility establishes that learning remains useful;
malicious ASR measures attack amplification; artifact poisoning shows that the
unsafe rule entered persistent storage; contamination and carryover harm show
that the rule survives after the original attack disappears.

% Main Table 1: do not place SafeEvolve in this table.
\begin{table*}[t]
\centering
\small
\setlength{\tabcolsep}{4.2pt}
\renewcommand{\arraystretch}{1.12}
\resizebox{\textwidth}{!}{
\begin{tabular}{lcc|cc|cc|ccc}
\toprule
& \multicolumn{2}{c|}{\textbf{Capability ($\uparrow$)}}
& \multicolumn{2}{c|}{\textbf{Attack amplification ($\downarrow$)}}
& \multicolumn{2}{c|}{\textbf{Evolved artifact ($\downarrow$)}}
& \multicolumn{3}{c}{\textbf{Post-attack behavior}} \\
\cmidrule(lr){2-3}\cmidrule(lr){4-5}\cmidrule(lr){6-7}\cmidrule(lr){8-10}
\textbf{Evolution setting}
& B-Util. & Carry-Util.
& M-ASR & $\Delta$ASR
& Poisoned & Gen. score
& Contam. $\downarrow$ & Carry-Harm $\downarrow$ & Onset $\downarrow$ \\
\midrule
No Evolution          & -- & -- & -- & 0.0 & 0.0 & 0.0 & -- & 0.0 & -- \\
Read-only Reflection  & -- & -- & -- & --  & --  & --  & -- & --  & -- \\
Benign Evolution      & -- & -- & -- & --  & --  & --  & -- & --  & -- \\
\midrule
\skillrevise          & -- & -- & -- & --  & --  & --  & -- & --  & -- \\
EvoSkill              & -- & -- & -- & --  & --  & --  & -- & --  & -- \\
SkillClaw             & -- & -- & -- & --  & --  & --  & -- & --  & -- \\
AutoSkill             & -- & -- & -- & --  & --  & --  & -- & --  & -- \\
SkillsVote            & -- & -- & -- & --  & --  & --  & -- & --  & -- \\
Hermes-native         & -- & -- & -- & --  & --  & --  & -- & --  & -- \\
\bottomrule
\end{tabular}}
\caption{\textbf{Skill evolution can preserve utility while creating
persistent unsafe behavior (RQ1).} Pooled results over the frozen episode
distribution. B-Util. is utility on benign interleaved tasks; Carry-Util. is
utility in the fresh-session persistence probe; M-ASR is attack success on
exposed tasks; $\Delta$ASR is relative to matched no evolution; Poisoned is the
rate of stored skill bundles judged unsafe and over-generalized; Gen. score is
artifact-level unsafe generalization; Contam. is unsafe-goal completion on
zero-injection benign tasks; Carry-Harm is unsafe-goal completion in the fresh
session; Onset is the mean number of exposures before first contamination.
Report 95\% episode-bootstrap intervals in parentheses or in an accompanying
interval plot. Best utility and lowest risk among evolution methods are bold.}
\label{tab:misevolution-main}
\end{table*}

\paragraph{Result paragraph template.}
Across evolution methods, benign utility changes by [X] while malicious ASR
increases by [Y]. More importantly, [Z]\% of episodes exhibit zero-injection
contamination and [W]\% retain harm in a fresh session. The joint result rules
out a pure context effect: the attack is absent, the executor is reset, and the
only cross-task object is the evolved library.

\paragraph{Safety--utility frontier.}
Plot each method by benign utility and carryover harm. The important region is
high utility and high harm: these methods learn useful procedures while also
retaining the unsafe shortcut, making capability improvement an unreliable
safety signal.

\subsection{RQ2: How does exposure become persistent harm?}
\label{sec:rq2}

RQ2 tests each link of Equation~\ref{eq:chain}. First report the probability
that an exposure produces a poisoned artifact. Then report retrieval of that
artifact on later benign tasks, contamination conditional on retrieval, and
fresh-session harm conditional on final-library poisoning. Use matched
mediation or stratified conditional effects rather than inferring mediation
from a single example.

\begin{table}[t]
\centering
\small
\setlength{\tabcolsep}{5pt}
\begin{tabular}{lcc}
\toprule
\textbf{Link in the chain} & \textbf{Rate} & \textbf{Conditional effect} \\
\midrule
Exposure $\rightarrow$ poisoned skill & -- & -- \\
Poisoned skill $\rightarrow$ retrieval & -- & -- \\
Retrieval $\rightarrow$ contamination & -- & -- \\
Final poison $\rightarrow$ carryover harm & -- & -- \\
\bottomrule
\end{tabular}
\caption{\textbf{Exposure-to-persistence chain (RQ2).} Each row measures an
observable transition rather than treating later harm as automatic evidence of
skill mediation. Conditional effects compare matched episodes with and without
the upstream event.}
\label{tab:chain}
\end{table}

\paragraph{Artifact risk precedes behavioral harm.}
Compare artifact poisoning with realized contamination. Episodes with poisoned
skills but no immediate behavioral harm constitute latent risk: the persistent
policy has changed, but the evaluated task has not supplied the retrieval cue
or execution opportunity. Quantify this artifact--behavior gap by method and
agent.

\paragraph{Worked causal trace.}
Show one compact case from source exposure to authored skill diff, later
retrieval log, unsafe tool action, judge state, and fresh-session reproduction.
The example should expose the exact over-generalized clause and the safety check
that evolution removed or failed to encode.

\subsection{RQ3: Where does misevolution emerge, and at what budget?}
\label{sec:rq3}

Report heterogeneity across evolution algorithms, agent frameworks, model
backbones, substrates, domains, and risk families. Avoid treating the pooled
mean as universality: identify which design properties predict artifact
poisoning and which executor policies block poisoned artifacts from becoming
behavior.

\begin{figure}[t]
  \centering
  \fbox{\parbox[c][3.7cm][c]{0.94\linewidth}{
    \centering
    \textbf{Figure 3 placeholder: budget--response curve.}\\[4pt]
    X-axis: adversarial exposures $k\in\{0,1,2,3\}$.\\
    Y-axes or panels: artifact poisoning, zero-injection contamination, and
    fresh-session carryover harm. Show method-level curves with bootstrap CIs.
  }}
  \caption{\textbf{Misevolution under bounded exposure (RQ3).} The minimum
  budget that produces persistent harm defines the practical attacker cost.}
  \label{fig:budget}
\end{figure}

\paragraph{Failure mechanisms.}
Cluster audited skill diffs into a small post-hoc taxonomy such as verification
skipping, authorization assumption, policy erosion, and unsafe
over-generalization. Report multi-label frequencies and annotator agreement.
Use the taxonomy to explain variation, not as the primary proof that
misevolution exists.

\subsection{RQ4: Can SafeEvolve preserve learning without preserving harm?}
\label{sec:rq4}

The mitigation evaluation uses held-out domains and episodes that are hidden
during evolution, repair-policy development, and threshold selection. Libraries
are frozen before evaluation. The primary evidence is behavioral: a mitigation
must reduce contamination and fresh-session harm while retaining the utility of
raw evolution. Artifact scores diagnose how it achieves that result but are not
the sole success criterion.

% Main Table 2: held-out mitigation table.
\begin{table*}[t]
\centering
\small
\setlength{\tabcolsep}{4.2pt}
\renewcommand{\arraystretch}{1.10}
\resizebox{\textwidth}{!}{
\begin{tabular}{ll|cccc|ccc|c}
\toprule
& & \multicolumn{4}{c|}{\textbf{Held-out safety ($\downarrow$)}}
& \multicolumn{3}{c|}{\textbf{Capability ($\uparrow$)}}
& \textbf{Cost} \\
\cmidrule(lr){3-6}\cmidrule(lr){7-9}\cmidrule(lr){10-10}
\textbf{Evolution method} & \textbf{Governance}
& $\Delta$Unsafe & Poisoned & Contam. & Carry-Harm
& B-Util. & Carry-Util. & Retained gain
& Overhead \\
\midrule
No Evolution & None & 0.0 & 0.0 & -- & 0.0 & -- & -- & 0.0\% & 1.00$\times$ \\
\midrule
\multirow{3}{*}{\skillrevise}
 & Raw             & -- & -- & -- & -- & -- & -- & 100\% & 1.00$\times$ \\
 & Library-Drift   & -- & -- & -- & -- & -- & -- & --    & -- \\
 & \safeevolve     & -- & -- & -- & -- & -- & -- & --    & -- \\
\midrule
\multirow{3}{*}{EvoSkill}
 & Raw             & -- & -- & -- & -- & -- & -- & 100\% & 1.00$\times$ \\
 & Library-Drift   & -- & -- & -- & -- & -- & -- & --    & -- \\
 & \safeevolve     & -- & -- & -- & -- & -- & -- & --    & -- \\
\midrule
\multirow{3}{*}{SkillClaw}
 & Raw             & -- & -- & -- & -- & -- & -- & 100\% & 1.00$\times$ \\
 & Library-Drift   & -- & -- & -- & -- & -- & -- & --    & -- \\
 & \safeevolve     & -- & -- & -- & -- & -- & -- & --    & -- \\
\midrule
\multicolumn{2}{l|}{\textbf{Macro-average over supported methods}}
 & -- & -- & -- & -- & -- & -- & -- & -- \\
\bottomrule
\end{tabular}}
\caption{\textbf{SafeEvolve reduces evolution-induced harm while preserving
the capability gain (RQ4).} Each evolution method is evaluated raw, with a
capability-oriented library hygiene baseline, and with \safeevolve on the same
held-out episodes. $\Delta$Unsafe aggregates paired behavior-level unsafe
change relative to no evolution; Poisoned is reported diagnostically; Retained
gain is $(U_{\mathrm{mitigated}}-U_{\mathrm{no\mbox{-}evo}})/
(U_{\mathrm{raw}}-U_{\mathrm{no\mbox{-}evo}})$. Overhead reports end-to-end
latency or model-call cost relative to raw evolution. Bold the safest method
whose utility is statistically non-inferior to raw evolution.}
\label{tab:safeevolve-main}
\end{table*}

\paragraph{Why Library-Drift is the decisive baseline.}
Capability-oriented hygiene can remove stale or low-value skills, but it should
retain a skill that improves utility through an unsafe shortcut. A gap between
Library-Drift and \safeevolve therefore tests whether an explicitly safety-aware
objective is necessary.

\paragraph{Component evaluation.}
Evaluate the admission critic independently on malicious and benign skill
corpora, reporting recall, AUROC, and benign false-positive rate. Evaluate
verification preservation, provenance rollback, and safety-aware retirement in
ablations. These component results explain the mitigation table; they do not
replace its held-out behavioral evidence.

\begin{table}[t]
\centering
\small
\setlength{\tabcolsep}{4pt}
\begin{tabular}{lccc}
\toprule
\textbf{SafeEvolve variant} & \textbf{Recall $\uparrow$} &
\textbf{Benign FPR $\downarrow$} & \textbf{Carry-Harm $\downarrow$} \\
\midrule
Full \safeevolve             & -- & -- & -- \\
$-$ admission critic         & -- & -- & -- \\
$-$ verification preservation& -- & -- & -- \\
$-$ provenance               & -- & -- & -- \\
$-$ safety-aware retirement  & -- & -- & -- \\
\bottomrule
\end{tabular}
\caption{\textbf{SafeEvolve component analysis.} Detection metrics come from
independent labeled skill corpora; carryover harm comes from the held-out
episode benchmark.}
\label{tab:safeevolve-ablation}
\end{table}

\section{Discussion}
\label{sec:discussion}

\paragraph{Skills are security-relevant policy.}
The central implication is that a skill library should not be governed as a
passive document store. A stored skill can alter tool choice, ordering,
verification, and authorization across future tasks. Admission, provenance,
rollback, and safety monitoring should therefore be part of the learning
boundary.

\paragraph{Latent and realized misevolution.}
Artifact poisoning and behavioral harm are distinct operational states. A
poisoned skill that has not yet fired is latent risk, not a safe artifact. A
runtime policy can temporarily block execution while the persistent library
continues to carry the unsafe rule. Product evaluation should report both.

\paragraph{Capability and safety share the same shortcut.}
The difficult cases are not useless malicious skills. They are procedures that
increase success by removing friction that was serving as a safeguard. This is
why disabling evolution loses capability and capability-only library hygiene
does not solve the problem.

\paragraph{Scope and limitations.}
The study focuses on readable skill artifacts produced by CLI and coding agents
and on bounded direct exposure in executable task environments. Results support
claims about evolution-induced incremental risk within this setting. Future
work should test indirect exposure, multimodal skills, opaque learned memories,
longer deployment horizons, interacting libraries, and adaptive attackers.

\section{Conclusion}
\label{sec:conclusion}

Skill self-evolution changes the temporal structure of agent safety. An unsafe
action need not end with the trajectory that produced it: evolution can extract
the successful shortcut, store it as a reusable skill, and apply it later when
the original attack is gone. \bench makes this transition measurable from
exposure through clean-session persistence. \safeevolve shows how to govern the
persistent update while preserving useful learning. Safe self-improvement
requires evaluating not only whether an agent succeeds today, but also what
today's success teaches it to do tomorrow.

\bibliography{iclr2026_conference}
\bibliographystyle{iclr2026_conference}

\appendix

\section{Benchmark Construction}
\label{app:construction}

Document episode generation, source-task sampling, benign twins, adversarial
exposures, risk-family alignment, deduplication, leakage checks, and held-out
splits. Include the complete dataset card and licenses.

\section{Agent and Method Configurations}
\label{app:configurations}

Report every agent framework, model endpoint, evolution implementation,
retrieval configuration, update cadence, prompt, decoding parameter, repository
commit, and applicability restriction. Separate unsupported combinations from
failed runs.

\section{Judges and Artifact Auditing}
\label{app:judges}

Provide the original substrate judges, artifact-unsafety rubric,
over-generalization rubric, auditor prompts, calibration examples, agreement
statistics, and adjudication procedure. State which models are disjoint from
the agent and evolution model families.

\section{Additional Results}
\label{app:additional-results}

Include per-agent, per-backbone, per-domain, per-risk-family, per-position, and
per-budget tables with confidence intervals. Report invalid episodes, runtime
failures, missing combinations, and sensitivity to aggregation choices.

\section{SafeEvolve Implementation}
\label{app:safeevolve-implementation}

Give pseudocode and prompts for admission, repair, verification preservation,
lineage tracking, safety-aware selection, quarantine, and retirement. Include
cost accounting and the policy for re-auditing repaired bundles.

\section{Case Studies}
\label{app:case-studies}

For each major failure family, show the source trajectory, skill diff, lineage,
retrieval event, later tool trajectory, judge evidence, and effect of
\safeevolve. Include negative cases where the artifact is poisoned but the
executor blocks it, and where evolution remains safe.

\end{document}
